\ifdefined\gkver\else\def\gkver{student}\fi
\documentclass[\gkver]{gaokaozhenti}
\begin{document}

\chapter{2026年广西}

\begin{enumerate}
\item
%% number: 1
%% paperName: 2026年广西高考第 1 题 4 分
%% typeId: 1
%% type: 单选题
%% chapter: 原子和原子核
%% point: 原子核的组成
%% method: 无
%% score: 4
%% degree: 950
%% duplicateId: 0
%% body:
我国科学家首次合成的新核素 $\ce{^{235}_{97}Bk}$ 与核素 $\ce{^{235}_{92}U}$ 具有相同的 \xzanswer{B}

\fourchoices[answer=B]
{核电荷数}
{质量数}
{原子质量}
{中子数}

%% answer: B

%% memo:
\memoanswer{
原子符号左上角的数值代表质量数，根据题意可知核素 $\ce{^{235}_{97}Bk}$ 与核素 $\ce{^{235}_{92}U}$ 质量数相等，质子数（核电荷数）不同；

根据“中子数 = 质量数 $-$ 质子数”，二者质子数不同、质量数相同，因此中子数也不同；

原子质量由原子核内所有质子、中子的质量和核外电子质量共同决定，二者质子、中子总数相同但质子和中子的各自数量不同，因此原子质量不相等。

故选 B。
}
\item
%% number: 2
%% paperName: 2026年广西高考第 2 题 4 分
%% typeId: 1
%% type: 单选题
%% chapter: 机械波
%% point: 波长频率波速
%% method: 无
%% score: 4
%% degree: 950
%% duplicateId: 0
%% body:
二胡发出的乐音中含频率为 $1000 \UHz$ 的声波，空气中的声速约为 $340 \Ums$，则该频率声波在空气中的波长约为 \xzanswer{A}

\fourchoices[answer=A]
{$0.34 \Um$}
{$2.94 \Um$}
{$3.40 \Um$}
{$340 \Um$}

%% answer: A

%% memo:
\memoanswer{
波速、波长、频率满足波的基本关系 $v = \lambda f$，

变形可得波长计算公式 $\lambda=\frac{v}{f}$，

代入题干数据 $v = 340 \Ums$、$f = 1000 \UHz$，

计算得 $\lambda=\frac{340}{1000}\Um=0.34\Um$。

故选 A。
}
\item
%% number: 3
%% paperName: 2026年广西高考第 3 题 4 分
%% typeId: 1
%% type: 单选题
%% chapter: 直线运动
%% point: 质点
%% method: 无
%% score: 4
%% degree: 950
%% duplicateId: 0
%% body:
若在平陆运河试运行时，要分析航道内货轮的运动，可忽略货轮形状和大小的是 \xzanswer{C}

\fourchoices[answer=C]
{研究货轮的调头过程}
{计算货轮穿过跨河大桥的时间}
{计算货轮航行 $10 \Ukm$ 的平均速率}
{研究货轮避让障碍物时的相对位置}

%% answer: C

%% memo:
\memoanswer{
质点的判定依据：当物体的大小、形状对所研究问题的影响可以忽略时，可将物体视为质点。

A．货轮调头时，研究的是船的姿态转动，船的形状、大小对该问题的影响不可忽略，故 A 错误；

B．货轮穿过跨河大桥时，需考虑船的长度与大桥的匹配关系，船的大小形状不可忽略，故 B 错误；

C．研究船以 $10 \Ukm$ 的平均速率行驶时，船的大小相对于运河的整体运动尺度可忽略，对运动速率的研究无影响，故 C 正确；

D．躲避障碍物时，需考虑船的大小、各部位位置以避免碰触，大小形状不可忽略，故 D 错误。

故选 C。
}
\item
%% number: 4
%% paperName: 2026年广西高考第 4 题 4 分
%% typeId: 1
%% type: 单选题
%% chapter: 电场
%% point: 电场线
%% method: 无
%% score: 4
%% degree: 850
%% duplicateId: 0
%% body:
天气干燥，人体因与衣服摩擦而带上正电，当手靠近金属门把手时，两者间形成静电场，$P$、$Q$ 是同一条电场线上的两点，如图，则 \xzanswer{D}

\onepicture[w=3.4cm]{\includesvg{figs/04}}

\fourchoices[answer=D]
{$P$ 点电势高于 $Q$ 点电势}
{$P$ 点电场强度大小大于 $Q$ 点电场强度大小}
{试探电荷 $+q$ 在 $P$ 点电势能高于其在 $Q$ 点电势能}
{手触碰门把手瞬间，电子从门把手流向人体}

%% answer: D

%% memo:
\memoanswer{
A．由于人手带正电，在门把手表面感应出负电荷，电场线由正电荷指向负电荷，故电场线由 $Q$ 指向 $P$，沿电场线方向电势逐渐降低，有 $\varphi_{Q} > \varphi_{P}$，故 A 错误；

B．$Q$ 点更靠近作为场源的手指，因此 $Q$ 点电场强度大于 $P$ 点，即 $E_{Q} > E_{P}$，故 B 错误；

C．正试探电荷的电势能 $E_{p} = q\varphi$，结合 $\varphi_{Q} > \varphi_{P}$，可知 $E_{pQ} > E_{pP}$，故 C 错误；

D．人手与门把手触碰时，电子从门把手流向人手，与正电荷进行中和，故 D 正确。

故选 D。
}
\item
%% number: 5
%% paperName: 2026年广西高考第 5 题 4 分
%% typeId: 1
%% type: 单选题
%% chapter: 运动和力
%% point: 动量定理
%% method: 无
%% score: 4
%% degree: 850
%% duplicateId: 0
%% body:
立定跳远的起跳过程可简化为人从半蹲姿势发力至双脚离地的过程。某同学训练后在起跳过程中对水平地面的平均作用力大小增大。若训练前后该同学起跳过程持续时间不变，起跳方向与地面夹角不变。则与训练前相比，该同学 \xzanswer{B}

\fourchoices[answer=B]
{离地瞬间的动量大小变小}
{离地瞬间的速度大小变大}
{起跳过程中受到地面的冲量不变}
{起跳过程中受到地面的平均支持力不变}

%% answer: B

%% memo:
\memoanswer{
ABC．由于起跳方向与作用时间 $\Delta t$ 不变，当用更大的平均作用力起跳时，根据牛顿第三定律可知，地面对人的平均作用力 $F$ 也增大，人所受的平均合力 $F_{\text{合}}$ 也增大，由动量定理有 $I = F_{\text{合}} t = \Delta p = m \Delta v$，

可知，起跳完成时动量变大，地面给人的冲量变大，由于初速度为 0，起跳完成时速度变大，故 AC 错误，B 正确；

D．由于起跳方向不变，设起跳方向与地面的夹角为 $\theta$，人的质量为 $m$，重力加速度为 $g$，则地面对人的平均支持力为 $F_{N} = F_{\text{合}} \sin \theta + mg$，

则地面的平均支持力变大，故 D 错误。

故选 B。
}
\item
%% number: 6
%% paperName: 2026年广西高考第 6 题 4 分
%% typeId: 1
%% type: 单选题
%% chapter: 磁场
%% point: 带电粒子在磁场中的运动
%% method: 无
%% score: 4
%% degree: 450
%% duplicateId: 0
%% body:
如图，在以坐标原点 $O$ 为圆心、半径为 $R$ 的圆形区域内，存在磁感应强度大小为 $B$、方向垂直纸面向里的匀强磁场。一带电粒子从 $M$ 点沿 $x$ 轴正方向射入，从 $N$ 点射出，出射方向与 $x$ 轴正方向夹角为 $60{^\circ}$。已知粒子电荷量为 $q$，质量为 $m$，不计重力。该粒子 \xzanswer{C}

\onepicture[w=4.6cm]{figs/06.png}

\fourchoices[answer=C]
{带正电荷}
{从 $M$ 到 $N$ 的运动轨迹是抛物线}
{射入磁场时的速率为 $\frac{\sqrt{3}qBR}{m}$}
{从 $M$ 到 $N$ 的运动时间 $\frac{2\pi m}{3qB}$}

%% answer: C

%% memo:
\memoanswer{
A．根据题意，由左手定则可知，带电粒子带负电，故 A 错误；

B．带电粒子在磁场中做匀速圆周运动，则粒子的运动轨迹 $MN$ 是圆弧，故 B 错误；

CD．根据题意，画出粒子的运动轨迹，如图所示

\onepicture[w=4.6cm]{figs/06_an.jpg}

设粒子做匀速圆周运动的半径为 $r$，由几何关系有 $\frac{r}{R} = \tan 60^\circ$，

解得 $r = \sqrt{3}R$。

由洛伦兹力提供向心力有 $qv_{0}B = m\frac{v_{0}^{2}}{r}$，

解得 $v_{0} = \frac{qBr}{m} = \frac{\sqrt{3}qBR}{m}$。

由几何关系可知，粒子在磁场中运动轨迹对应的圆心角为 $\theta = 60^\circ$，则粒子在磁场内的运动时间 $t = \frac{\theta}{2\pi} \cdot \frac{2\pi r}{v_0} = \frac{\pi m}{3qB}$，故 C 正确，D 错误。

故选 C。
}
\item
%% number: 7
%% paperName: 2026年广西高考第 7 题 4 分
%% typeId: 1
%% type: 单选题
%% chapter: 电磁感应
%% point: 切割磁感线电动势
%% method: 无
%% score: 4
%% degree: 450
%% duplicateId: 0
%% body:
某发光装置的电路如图~\ref{2026广西07a}。金属棒 $OP$ 一端固定在竖直导电转轴 $OO'$ 的 $O$ 点，另一端固定在水平导电圆环的 $P$ 点，圆环圆心为 $O$、半径为 $l$。圆环和转轴分别通过电刷 $M$、$N$ 保持与外电路的连接。圆环内存在磁感应强度大小为 $B$、方向竖直向下的匀强磁场。转动转轴时，金属棒 $OP$ 和圆环随转轴一起以角速度逆时针转动。$t = 0$ 时刻，发光二极管 $LED_{1}$ 和 $LED_{2}$ 亮度相同，$\omega$ 随 $t$ 的变化关系如图~\ref{2026广西07b}。已知电路中金属棒 $OP$、电阻 $R$、电感线圈 $L$ 的电阻，$LED_{1}$ 和 $LED_{2}$ 正向导通的电阻均为 $r$，其余电阻不计。则 \xzanswer{B}

\twopicture[labelA=2026广西07a, labelB=2026广西07b, h=3.4cm]
{figs/07a.png}
{figs/07b.png}

\fourchoices[answer=B]
{$t = 0$ 到 $t = \frac{2\pi}{\omega_{0}}$ 的过程中，通过 $LED_{1}$ 的电流为 $\frac{Bl^{2}\omega_{0}}{4r}$}
{$t = 0$ 到 $t = \frac{2\pi}{\omega_{0}}$ 的过程中，整个装置消耗的电能为 $\frac{\pi B^{2}l^{4}\omega_{0}}{4r}$}
{$t = t_{1}$ 时，电感线圈 $L$ 两端电压为 0}
{$t = \frac{2\pi}{\omega_{0}}$ 到 $t = t_{2}$ 的过程中，$LED_{1}$ 比 $LED_{2}$ 先熄灭}

%% answer: B

%% memo:
\memoanswer{
A．初始时，金属棒 $OP$ 转动切割磁感线，感应电动势为 $E_{1}=\frac{1}{2}Bl^{2}\omega_{0}$，

电路总电阻为 $R_{\text{总}}=\frac{2r\cdot2r}{2r+2r}+r=2r$，

则金属棒 $OP$ 两端的电压为

\[ U=E_{1}-\frac{E_{1}}{R_{\text{总}}}\cdot r=\frac{1}{2}E_{1}=\frac{1}{4}Bl^{2}\omega_{0} \]

通过 $LED_{1}$ 的电流 $I_{1}=\frac{U}{2r}=\frac{Bl^{2}\omega_{0}}{8r}$，故 A 错误；

B．$0 \sim t_{0}$ 时间内，整个电路消耗的电能 $E_{\text{电}}=\frac{E_{1}^{2}}{R_{\text{总}}}\cdot t_{0}$，

其中 $t_{0}=\frac{2\pi}{\omega_{0}}$，得 $E_{\text{电}}=\frac{\pi B^{2}l^{4}\omega_{0}}{4r}$，故 B 正确；

C．$t_{0} \sim t_{1}$ 时间内角速度在变化，金属棒 $OP$ 产生的感应电动势在变化，电路中的电流在变化，由于电感 $L$ 的自感作用，会产生自感电动势阻碍电流的变化，所以 $t_{1}$ 时刻电感 $L$ 两端电压不为零，故 C 错误；

D．$t_{1} \sim t_{2}$ 时间内，金属棒 $OP$ 的角速度为 0，产生的感应电动势为 0，电感 $L$ 会产生自感电动势，自感电流与原电流方向相同，$LED_{1}$ 与电感 $L$ 串联，$LED_{2}$ 与电阻 $R$ 串联，自感电流无法反向通过 $LED_{2}$，$\frac{2\pi}{\omega_{0}} \sim t_{2}$ 时间，$LED_{1}$ 比 $LED_{2}$ 后熄灭，故 D 错误。

故选 B。
}
\item
%% number: 8
%% paperName: 2026年广西高考第 8 题 6 分
%% typeId: 2
%% type: 多选题
%% chapter: 曲线运动
%% point: 斜抛运动
%% method: 无
%% score: 6
%% degree: 650
%% duplicateId: 0
%% body:
除冰机器人清除输电线上的覆冰时，若破碎的冰块从同一点向四周飞溅，仅在重力作用下运动并落在水平地面上，则落地过程所用时间相等的是 \xzanswer{AD}

\fourchoices[answer=AD]
{水平飞出的冰块}
{质量相等的冰块}
{向下飞出的冰块}
{初速度相同的冰块}

%% answer: AD

%% memo:
\memoanswer{
A．若要碎冰块同时落地，则碎冰块在竖直方向的初速度需要相同，水平抛出，竖直方向的初速度为 0，满足条件，A 正确；

BC．质量相等和向下飞出的冰块不能保证竖直方向的初速度相同，BC 错误；

D．初速度相同（点拨：速度是矢量，既有大小又有方向），则竖直方向的初速度相同，满足条件，D 正确。

故选 AD。
}
\item
%% number: 9
%% paperName: 2026年广西高考第 9 题 6 分
%% typeId: 2
%% type: 多选题
%% chapter: 光学
%% point: 光的折射
%% method: 无
%% score: 6
%% degree: 450
%% duplicateId: 0
%% body:
两块楔形电光材料置于两平行金属电极之间，通过图~\ref{2026广西09a} 所示电路施加电压实现光束偏转，右视图如图~\ref{2026广西09b}。两电光材料的折射率与电极间的电压 $U$ 分别满足关系 $n_{1} = n_{0} + kU$ 和 $n_{2} = n_{0} - kU$，其中 $U > 0$ 且在安全电压范围内，$k > 0$，$n_{2} > 1$。空气中一束光从正前方沿 $OO'$ 垂直电光材料 1 的 $MN$ 界面入射，最终从电光材料 2 的 $PQ$ 界面出射，出射光线相对于入射光线偏转的角度为 $\theta$。已知所有界面均不发生全反射，且只考虑折射。初始时刻，滑动变阻器 $R$ 的滑片置于 $b$ 端，闭合开关 $S$。则 \xzanswer{BD}

\twopicture[labelA=2026广西09a, labelB=2026广西09b, h=3.4cm]
{figs/09a.png}
{figs/09b.png}

\fourchoices[answer=BD]
{光线从 $O'P$ 区域出射}
{光线从 $O'Q$ 区域射出}
{滑动变阻器的滑片从 $b$ 向 $a$ 移动时，$\theta$ 逐渐增大}
{滑动变阻器的滑片从 $b$ 向 $a$ 移动时，$\theta$ 逐渐减小}

%% answer: BD

%% memo:
\memoanswer{
AB．由于 $U > 0$，$k > 0$ 始终有 $n_{1} > n_{2}$。

设界面与 $PQ$ 边夹角为 $\alpha$，光束在两材料界面折射时，折射角为 $r$，则入射角 $i = \alpha$。

根据折射定律有 $n_{1} \sin i = n_{2} \sin r$。

由于 $n_{1} > n_{2}$，可知 $r > i$。

法线垂直于界面，由几何关系可知，入射光线水平向右，折射光线向下偏折（点拨：光由光密介质射入光疏介质，折射角大于入射角，光线向界面靠拢，即向下偏），作出光路图如图所示。

\onepicture[w=4.8cm]{figs/09_an.jpg}

故折射光从 $O$ 下方的 $O'Q$ 区域射出，故 A 错误，B 正确；

CD．由电路图可知，两种材料并联在滑动变阻器滑片左侧部分，则当滑动变阻器的滑片向左移动时，材料两端电压 $U$ 减小，根据题意可知 $n_{1}$ 减小、$n_{2}$ 增大，所以 $\sin r$ 减小，折射角 $r$ 减小，折射光线与水平面的夹角 $\theta = r - i$，所以 $\theta$ 变小，故 C 错误，D 正确。

故选 BD。
}
\item
%% number: 10
%% paperName: 2026年广西高考第 10 题 6 分
%% typeId: 2
%% type: 多选题
%% chapter: 曲线运动
%% point: 人造地球卫星
%% method: 无
%% score: 6
%% degree: 350
%% duplicateId: 0
%% body:
如图，两颗人造卫星的轨道可视为以地心 $O$ 为圆心的共面同心圆，卫星 1、卫星 2 的轨道半径分别为 $R_{1}$、$R_{2}$（$R_{2} > R_{1}$）。初始时刻，卫星 1 位于 $P$ 点、卫星 2 位于 $Q$ 点，此时卫星 1 沿自身轨道切线方向发射探测器，探测器沿椭圆轨道无动力飞行，第一次到达 $K$ 点时恰好与卫星 2 相遇，且 $P$、$O$、$K$ 三点共线。已知地球质量为 $M$，引力常量为 $G$，图中箭头所示为卫星及探测器的绕行方向。若卫星发射探测器时自身的速度变化忽略不计，则 \xzanswer{AC}

\onepicture[w=4.8cm]{figs/10.png}

\fourchoices[answer=AC]
{卫星 2 的运行速度大小 $\sqrt{\frac{GM}{R_{2}}}$}
{探测器沿椭圆轨道运动的周期 $2\pi\sqrt{\frac{(R_{1} + R_{2})^{3}}{GM}}$}
{图中 $\angle POQ$ 等于 $\pi - \frac{\pi}{4}\sqrt{\frac{2(R_{1} + R_{2})^{3}}{R_{2}^{3}}}$}
{两卫星从初始时刻到第一次相距最近时，历时 $\frac{2\pi\sqrt{R_{1}^{3}R_{2}^{3}}-\pi\sqrt{\frac{2R_{1}^{3}(R_{1}+R_{2})^{3}}{R_{2}^{3}}}}{2\sqrt{GM}\left(\sqrt{R_{2}^{3}}-\sqrt{R_{1}^{3}}\right)}$}

%% answer: AC

%% memo:
\memoanswer{
A．根据万有引力提供向心力有 $\frac{GMm}{R_{2}^{2}} = m\frac{v_{2}^{2}}{R_{2}}$，解得 $v_{2} = \sqrt{\frac{GM}{R_{2}}}$，故 A 正确；

B．由几何关系可知，椭圆轨道的半长轴为 $a = \frac{R_{1} + R_{2}}{2}$，由开普勒第三定律有 $\frac{a^{3}}{T^{2}} = \frac{R_{2}^{3}}{T_{2}^{2}}$，

由万有引力提供向心力有 $\frac{GMm}{R_{2}^{2}} = m\frac{4\pi^{2}}{T_{2}^{2}}R_{2}$，

联立解得 $T = 2\pi\sqrt{\frac{(R_{1} + R_{2})^{3}}{8GM}}$，故 B 错误；

C．变轨卫星从 $P$ 到 $K$ 经过半个椭圆，则运动时间 $t = \frac{1}{2}T = \pi\sqrt{\frac{(R_{1} + R_{2})^{3}}{8GM}}$，

卫星在高轨道 2 的角速度为 $\omega_{2} = \frac{v_{2}}{R_{2}} = \sqrt{\frac{GM}{R_{2}^{3}}}$，

则高轨卫星从 $Q$ 到 $K$ 转过的角度 $\angle QOK = \omega_{2}t = \frac{\pi}{4}\sqrt{\frac{2(R_{1} + R_{2})^{3}}{R_{2}^{3}}}$，

则有 $\angle POQ = \pi - \frac{\pi}{4}\sqrt{\frac{2(R_{1} + R_{2})^{3}}{R_{2}^{3}}}$，故 C 正确；

D．由上述分析同理可得，卫星在轨道 1 的角速度为 $\omega_{1} = \sqrt{\frac{GM}{R_{1}^{3}}}$，

由于 $R_{2} > R_{1}$，则有 $\omega_{1} > \omega_{2}$，

由几何关系可知，第一次相距最近时，相对转过的角度等于初始位置角度差，即

$\left(\omega_{1} - \omega_{2}\right)t = \angle POQ = \pi - \frac{\pi}{4}\sqrt{\frac{2\left(R_{1} + R_{2}\right)^{3}}{R_{2}^{3}}}$，

解得 $t = \frac{\sqrt{R_{1}^{3}R_{2}^{3}}\left[\pi - \frac{1}{4}\pi\sqrt{\frac{2(R_{1} + R_{2})^{3}}{R_{2}^{3}}}\right]}{\sqrt{GM}\left(\sqrt{R_{2}^{3}} - \sqrt{R_{1}^{3}}\right)}$，故 D 错误。

故选 AC。
}
\item
%% number: 11
%% paperName: 2026年广西高考第 11 题 6 分
%% typeId: 4
%% type: 实验题
%% chapter: 力物体的平衡
%% point: 力
%% method: 无
%% score: 6
%% degree: 650
%% duplicateId: 0
%% body:
某同学按图~\ref{2026广西11a} 电路测量干电池的电动势和内阻，所用器材有电压表（量程 $0 \sim 3 \UV$，可视为理想电压表）、定值电阻 $R_{0} = 10.0 \UO$、开关 $S_{1}$ 和 $S_{2}$、干电池、导线若干。

\twopicture[labelA=2026广西11a, labelB=2026广西11b, h=3.2cm]
{figs/11a.png}
{figs/11b.png}

\begin{enumerate}
	\item
	断开 $S_{2}$，再闭合 $S_{1}$，电压表读数为 $1.44 \UV$，则干电池电源电动势的测量值 $E =$ \tkanswer[1.6cm]{1.44} $\UV$；
	\item
	保持 $S_{1}$ 闭合，再闭合 $S_{2}$，电压表示数 $U_{2}$，如图~\ref{2026广西11b}，则 $U_{2} =$ \tkanswer[1.6cm]{1.20} $\UV$，干电池内阻的测量值 $r =$ \tkanswer[1.6cm]{2.0} $\UO$。
\end{enumerate}

%% answer:

\jdanswer{
\begin{enumerate}
	\item
	$1.44$

	\item
	$1.20$，$2.0$

\end{enumerate}
}

%% memo:
\memoanswer{
（1）只闭合 $S_{1}$，理想电压表所在支路相当于断路，电压表示数为电源电动势，则电源电动势 $E = 1.44 \UV$。

（2）由题图~\ref{2026广西11b} 可知，电压表的分度值为 $0.1 \UV$，电压表读数为 $1.20 \UV$。

$S_{1}$、$S_{2}$ 均闭合时，由闭合电路欧姆定律得

\[ E=U+\frac{U}{R_{0}}r \]

代入数据解得 $r = 2.0 \UO$。
}
\item
%% number: 12
%% paperName: 2026年广西高考第 12 题 10 分
%% typeId: 4
%% type: 实验题
%% chapter: 气体性质
%% point: 验证玻意耳定律
%% method: 无
%% score: 10
%% degree: 450
%% duplicateId: 0
%% body:
某同学利用注射器、压强传感器、塑料管等器材，在 $25\UduC$ 的恒温环境下做“探究气体等温变化的规律”实验。

\twopicture[labelA=2026广西12a, labelB=2026广西12b, h=2.9cm]
{figs/12a.png}
{figs/12b.png}

\begin{enumerate}
	\item
	实验过程中，缓慢推动注射器活塞改变体积是为了保证封闭气体的 \tkanswer[1.8cm]{温度} 不变；
	\item
	由注射器壁上的刻度读出的是注射器内气体的 \tkanswer[1.8cm]{体积}，由压强传感器测得气体的压强；
	\item
	在测量某次数据时，该同学出现如图~\ref{2026广西12a} 所示的误操作，若体积读数准确且传感器工作正常，此时记录的压强值与 $25\UduC$ 时同体积下压强的理论值相比，会 \tkanswer[1.8cm]{偏大}；
	\item
	完成实验后，根据实验数据作 $\frac{1}{p}-V$ 图像，如图~\ref{2026广西12b}，得到一条直线，但不经过原点。请分析图~\ref{2026广西12b} 中直线不经过原点的原因并提出改进实验的措施（写出一种即可）。
	\begin{steps}
		\item
		原因：\tkanswer[6.0cm]{忽略了注射器与传感器连接的塑料管中残留气体的体积}；
		\item
		措施：\tkanswer[6.0cm]{减小塑料管的长度、提前测定残留体积并修正测量数据、改用内径更小的塑料管等}。
	\end{steps}
\end{enumerate}

%% answer:

\jdanswer{
\begin{enumerate}
	\item
	温度

	\item
	体积

	\item
	偏大

	\item
	①忽略了注射器与传感器连接的塑料管中残留气体的体积；②减小塑料管的长度、提前测定残留体积并修正测量数据、改用内径更小的塑料管等

\end{enumerate}
}

%% memo:
\memoanswer{
（1）缓慢推动活塞，注射器内的气体有充足时间和外界发生热交换，保持温度不变。

（2）注射器上的刻度标注的是内部空腔的容积，当封闭一定量气体在注射器内时，注射器对应的刻度就是封闭气体的体积。

（3）某同学手握注射器时，人体温度高于环境温度，热量传递给封闭气体，气体温度升高，由理想气体状态方程 $\frac{pV}{T}=C$ 可知，封闭气体的体积 $V$ 不变、热力学温度 $T$ 升高，则测量的气体压强 $p$ 大于理论压强。

（4）由 $\frac{1}{p} - V$ 图像可知，气体体积为 0 时，压强不为无穷大，说明注射器和传感器间有部分气体体积未被测量，即忽略了注射器与传感器连接的塑料管中残留气体的体积；改进措施是减小塑料管的长度、提前测定残留体积并修正测量数据、改用内径更小的塑料管等。
}
\item
%% number: 13
%% paperName: 2026年广西高考第 13 题 10 分
%% typeId: 6
%% type: 计算题
%% chapter: 运动和力
%% point: 牛顿第二定律
%% method: 无
%% score: 10
%% degree: 550
%% duplicateId: 0
%% body:
如图，工人使用夹具移动石料，夹具由不可伸长的绳 $AB$、$BC$、$BC'$ 及硬杆 $CD'$、$C'D$ 组成，两杆在 $O$ 点铰接。若长方体石料质量 $m = 80 \Ukg$，质量分布均匀，夹具夹起石料时，石料底面保持水平，$BCOC'$ 为正方形。忽略绳和杆的质量，重力加速度 $g$ 取 $10 \Umsq$。

\begin{enumerate}
	\item
	夹具夹着石料离开地面静止时，求绳 $BC$ 中的拉力大小；
	\item
	工人放下石料时，夹具夹着石料由静止开始匀加速下落 $0.1 \Um$，用时 $0.2 \Us$，求此过程中石料受到的摩擦力大小。
\end{enumerate}

\onepicture[w=6.0cm, align=right]{\includesvg{figs/13}}

%% answer:

\jdanswer{
\begin{enumerate}
	\item
	$400\sqrt{2}\UN$，沿 $BC$ 向下

	\item
	$400\UN$，方向竖直向上

\end{enumerate}
}

%% memo:
\memoanswer{
（1）以钩爪和石块组成的系统为研究对象，其受力分析如图所示。

\onepicture[w=3.0cm]{figs/13_an.jpg}

静止时系统受力平衡，由平衡条件有 $2F_{BC}\cos45^\circ=mg$，

代入数据解得 $F_{BC}=400\sqrt{2}\UN$。

由牛顿第三定律可得，静止时 $BC$ 绳受到钩爪的拉力

\[ F_{BC'} = F_{BC} = 400\sqrt{2}\UN \]

方向沿 $BC$ 向下。

（2）根据题意，由运动学公式有 $x=\frac{1}{2}at^{2}$，

代入数据解得 $a = 5 \Umsq$。

对石块由牛顿第二定律有 $mg - f = ma$，

代入数据解得 $f = 400\UN$，

方向竖直向上。
}
\item
%% number: 14
%% paperName: 2026年广西高考第 14 题 12 分
%% typeId: 6
%% type: 计算题
%% chapter: 功和能
%% point: 交通工具启动
%% method: 无
%% score: 12
%% degree: 450
%% duplicateId: 0
%% body:
某学习小组对一辆纯电动汽车的能耗开展研究：

\begin{enumerate}
	\item
	电池放电时能输出的总电荷量称作电池容量，通常以安培小时（$\UAh$）为单位，$1\UAh$ 等于多少库仑（$\UC$）电荷量？
	\item
	该车的电池容量为 $240 \UAh$，电池输出电压为 $375 \UV$，求该车电池充满电后可输出的电能 $E$。
	\item
	若该车在平直高速公路上行驶，受到的阻力来自轮胎与路面间的 $F_{1} = k_{1}Mg$ 和空气阻力 $F_{2} = k_{2}v^{2}$，其中 $k_{1}$、$k_{2}$ 为系数，$M$ 为车的质量，$v$ 为行驶速度，$g$ 为重力加速度，以上物理量的单位均采用国际单位制表示。该车电池输出电能转化为牵引力做功的效率为 $\eta$，求该车在平直高速公路上匀速行驶时，电池的输出功率 $P$ 的表达式。
	\item
	该车充满电后，计划 4 小时内在平直高速公路上行驶 400 千米。该车 $M = 2000 \Ukg$，$k_{1} = 0.0144$，$k_{2} = 0.36 \Ukgm$，当 $80 \Ukmh \leq v \leq 110 \Ukmh$ 时 $\eta = 0.9$，当 $v = 120 \Ukmh$ 时 $\eta = 0.8$；重力加速度 $g$ 取 $10 \Umsq$；请从时间和能耗的角度分析如下三种方案能否达成上述计划。

	\begin{tabular}{|c|l|}
	\hline
	方案 1 & 以 $90 \Ukmh$ 的速度匀速行驶 \\
	\hline
	方案 2 & 以 $108 \Ukmh$ 的速度匀速行驶 \\
	\hline
	方案 3 & 以 $120 \Ukmh$ 的速度匀速行驶 \\
	\hline
	\end{tabular}
\end{enumerate}

%% answer:

\jdanswer{
\begin{enumerate}
	\item
	$3600\UC$

	\item
	$3.24 \times 10^{8}\UJ$

	\item
	$\frac{(k_{1}Mg + k_{2}v^{2})v}{\eta}$

	\item
	方案二

\end{enumerate}
}

%% memo:
\memoanswer{
（1）由于 $1\Uh = 3600\Us$，所以 $q = It = 1\UA \times 3600\Us = 3600\UC$。

（2）因为 $q = It = 240\UA \times 3600\Us = 240 \times 3600\UC$，

所以根据题意有 $E = qU = 240 \times 3600 \times 375 \UJ = 3.24 \times 10^{8}\UJ$。

（3）根据题意可知汽车匀速行驶时，牵引力等于总阻力，即 $F = F_{1} + F_{2} = k_{1}Mg + k_{2}v^{2}$，

汽车牵引力的功率为 $P = Fv$，

又因为电池的输出功率 $P_{\text{出}} = \frac{P}{\eta} = \frac{(k_1Mg + k_2v^2)v}{\eta}$。

（4）根据题意行驶 400 千米，要求在时间 $t = 4\Uh$ 内完成，则全程最低平均速度

\[ \overline{v}_{\min} = \frac{s}{t} = \frac{400}{4} \Ukmh = 100 \Ukmh \]

由于方案一中 $90 \Ukmh < 100 \Ukmh$，故不满足时间要求；

方案二中 $v = 108 \Ukmh = 30 \Ums$，则由题可知 $\eta = 0.9$，此时有

\[ F = F_{1} + F_{2} = k_{1}Mg + k_{2}v^{2} = 612\UN \]

电池输出功率 $P_{\text{出}} = \frac{P}{\eta} = \frac{Fv}{\eta} = 20400\UW$，

整个过程行驶的时间 $t = \frac{400}{108}\Uh = \frac{100}{27}\Uh = \frac{100}{27} \times 3600\Us$，

所以消耗的能量为 $W = P_{\text{出}}t = 20400 \times \frac{100}{27} \times 3600\UJ = 2.72 \times 10^{8}\UJ$。

方案三中 $v = 120\Ukmh = \frac{100}{3}\Ums$，则由题可知 $\eta = 0.8$，此时有

\[ F' = F_{1} + F_{2} = k_{1}Mg + k_{2}v^{2} = 688\UN \]

电池输出功率 $P_{\text{出}}' = \frac{F'v}{\eta} = \frac{\frac{100}{3} \times 688}{0.8}\UW = \frac{344000}{12}\UW = \frac{86000}{3}\UW$，

整个过程行驶的时间 $t' = \frac{400}{120}\Uh = 12000\Us$，

所以消耗的能量为 $W' = P_{\text{出}}'t' = \frac{86000}{3} \times 12000\UJ = 3.44 \times 10^{8}\UJ$。

根据前面的计算可知电池的最大储能为 $E = 3.24 \times 10^{8}\UJ < 3.44 \times 10^{8}\UJ$，故方案三也不可行；所以应该选择方案二。
}
\item
%% number: 15
%% paperName: 2026年广西高考第 15 题 16 分
%% typeId: 6
%% type: 计算题
%% chapter: 磁场
%% point: 安培力
%% method: 无
%% score: 16
%% degree: 350
%% duplicateId: 0
%% body:
如图，水平桌面上固定着一对间距为 $l$ 的平行金属导轨。导轨间虚线包围区域由左侧的矩形和右侧的等腰三角形组成，等腰三角形底边为 $l$，底边上的高为 $h$。虚线包围的区域内存在磁感应强度大小为 $B$、方向垂直桌面向下的匀强磁场。一质量为 $m$ 的金属导体棒垂直放置在导轨上，初始位置位于矩形和等腰三角形连接处。设导体棒通有如图所示的电流 $I$，电流 $I$ 大小可调节，但每次运动过程中 $I$ 保持不变。导体棒与导轨间的动摩擦因数为 $\mu$（设大静摩擦力等于滑动摩擦力），重力加速度为 $g$。

\begin{enumerate}
	\item
	当 $I = I_{0}$ 时，求导体棒在初始位置受到的安培力大小和方向；
	\item
	对于 $I = I_{0}$ 情形，导体棒由初始位置从静止开始运动，求其能达到的最大速度的大小及其达到最大速度时位移的大小；
	\item
	若导体棒由初始位置从静止开始运动，停止时没有离开磁场区域，求 $I$ 的取值范围，并推导整个运动过程中安培力做的功与电流 $I$ 的函数关系。
\end{enumerate}

\onepicture[w=9.0cm, align=right]{\includesvg{figs/15}}

%% answer:

\jdanswer{
\begin{enumerate}
	\item
	$BI_{0}l$，垂直于棒向右

	\item
	$v = (BI_{0}l - \mu mg)\sqrt{\frac{h}{BI_{0}lm}}$，$x_{0} = h\left(1 - \frac{\mu mg}{BI_{0}l}\right)$

	\item
	$W_{\text{安}} = \frac{1}{2}BIlh$，$\frac{\mu mg}{Bl} < I \leq \frac{2\mu mg}{Bl}$

\end{enumerate}
}

%% memo:
\memoanswer{
（1）由题可知，在 $t=0$ 时，$F_{\text{安}}=BI_{0}l$，

由左手定则可知金属棒受到安培力方向垂直于棒向右。

（2）当金属棒达到最大速度时，此时安培力与摩擦力平衡，即 $BI_{0}L_{0} = \mu mg$，

设此时运动的位移为 $x_{0}$，根据几何关系可知 $\frac{h - x_{0}}{h} = \frac{L_{0}}{l}$，

解得 $L_{0}=\frac{\mu mg}{BI_{0}}$，$x_{0}=h\left(1-\frac{\mu mg}{BI_{0}l}\right)$。

在金属棒运动过程中只有安培力和摩擦力做功，有

\[ \sum BI_{0}L\cdot\Delta x-\mu mgx_{0}=\frac{1}{2}mv^{2} \]

即 $\frac{1}{2}BI_{0}\left(l+\frac{\mu mg}{BI_{0}}\right)\cdot h\left(1-\frac{\mu mg}{BI_{0}l}\right)-\mu mg\cdot h\left(1-\frac{\mu mg}{BI_{0}l}\right)=\frac{1}{2}mv^{2}$，

解得 $v = (BI_{0}l - \mu mg) \sqrt{\frac{h}{BI_{0}lm}}$。

（3）若金属棒恰好能到达 $O$ 点，安培力做功与电流的关系为 $W_{\text{安}} = \frac{1}{2} BIlh$，

根据动能定理，电流 $I$ 的最大值满足 $\frac{1}{2}BI_{\max}l\cdot h-\mu mg\cdot h=0$，

解得 $I_{\max} = \frac{2\mu mg}{Bl}$。

使金属棒恰好静止时电流满足 $BI_{\min}l = \mu mg$，

解得 $I_{\min} = \frac{\mu mg}{Bl}$。

为了使金属棒能运动起来且不离开磁场区域，则电流 $I$ 的取值范围为 $\frac{\mu mg}{Bl} < I \le \frac{2\mu mg}{Bl}$。
}
\end{enumerate}

\end{document}
